% Guide paragraphs written incrementally after image collation of each section.
\paragraph{110. Schematization through Postnikov fibrations, and its obstruction (pp. 443--445; 134-7/1--3).}
Grothendieck resumes the question of homotopy types over a ground ring by comparing an abelian group $M$ with the functor $W(M)(k)=M\otimes_{\underline Z}k$: the discrete object acquires a structure analogous to a vector bundle, suggesting an analogous construction for semisimplicial models and successive Postnikov fibers. The proposed ``schematization above level n'' would leave a lower truncation discrete and require the comparison from schematic cohomology with quasi-coherent coefficients to ordinary cohomology to be bijective. The reference to section 94 situates the hope of replacing $\underline Z$ by $\underline Q$ when nilpotence replaces an abelian fundamental group. On p.~444 he immediately questions this construction: torsion and the universal coefficients formula prevent the homotopy groups after base change from being obtained by tensoring the individual homotopy groups. The alternative is to reconstruct a homotopy type from a suitably enriched ``abelianization'' or ``linearization''; section 111 then seeks a different construction without explicit Postnikov decomposition. The handwritten replacement of ``(= homotopy)'' by ``(or rather, homotopy)'' on p.~444 matters here: the arXiv witness retains the cancelled equality wording. Its mapsto in the displayed base-change formula also differs from the scan's ordinary arrow; the scan alternates star subscripts with the typed notation X. Page 445 closes with a methodological reflection on the tension between detailed preparatory work and the desire to move quickly. The rust-obscured words on p.~443 remain marked rather than restored from the mathematical argument.

\paragraph{111. Unipotent bundles, divided powers, and a revised notion of equivalence (pp. 446--459; 134-7/4--17).}
The starting problem is to construct categories $Hot_n(k)$ of schematized $n$-connected homotopy types, with extension and restriction of the ground ring and a comparison with the ordinary homotopy category over $\underline Z$ (pp.~446--449). Grothendieck builds $U(k)$ inductively from affine bundles $W(M)$ and torsors, first in the fpqc setting. He asks whether the successive fibration data belong to the object or merely witness its existence, then proposes semisimplicial objects subject to connectivity and filtration conditions. The proposed weak equivalences are initially detected by taking sections, and the hoped-for $k$-module structures on homotopy groups lead to a search for a linearization into $D.(Ab_k)$. On p.~448 the model-category symbol has a typed subscript $m$, although the surrounding development uses $n$; this apparent index slip is retained in the text.

Pages 450--455 test divided powers as the required linearization. The candidate universal map $x\mapsto\sum x^{(i)}$ immediately forces a completion, and the completed divided-power coalgebra is considered through its group-like elements. After the Eilenberg--Maclane test, the discussion returns to pointed objects and uncompleted divided powers. The comparison (w) on p.~454 places ordinary free abelian linearization, $L$, and the pointed version $L_{pt}$ in one diagram. The marginal correction on p.~452 is consequential: ``One has to take k=Z, and components $X_n$ projective.'' Both its initial words and the subscript are legible on the scan. The subsequent distinction between quasi-isomorphisms and chain homotopisms supplies the projectivity restriction. The final reference in the sentence on p.~455 is to (x''), after a preceding reference to (x'''); the arXiv witness repeats (x''') and loses this distinction.

The argument is substantially revised on pp.~456--459. The author retracts the unrestricted derived-functor claim of p.~454, requires projective or flat components, and remarks that the Postnikov approach must again be set aside. The plain typed $I_k(M_*')$ on p.~456 is retained as printed, although the neighboring formulas suggest a notation slip. Page 457 then exposes the problem with ordinary homology as a $k$-module: the exterior-power calculation for $K(M,1)$, applied to $k=\underline Z[T]/(T^2+1)$, yields an obstruction. Page 458 accordingly defines total linearized homology $LH.(X_*)=L_{pt}(X_*)$ and considers weak equivalences detected by its homology invariants. The handwritten revision says it is becoming \emph{dubious} whether equivalences should be defined through $X_*(k)$; the arXiv witness reads \emph{obvious}, reversing the author's change of judgment. At the end of the section, Question 2 asks whether the two definitions coincide for projective pointed unipotent bundles satisfying $X_0=X_1=e$ (p.~459). The weaker condition $X_0=e$ raises an additional issue about abelian fundamental groups. These pages formulate questions and revised expectations, rather than establish an unrestricted equivalence theorem; section 112 resumes from the projective $\underline Z$ case.

\paragraph{112. Kan conditions and the scope of the linearization program (pp. 460--465; 134-7/18--23).}
The entry dated 26.8 takes the projective $\underline Z$ case of the comparison maps (w) as its working foothold. Grothendieck sketches a passage from the Eilenberg--Maclane case to finite chain complexes, then through finite and ``way-out'' infinite successive fibrations, using multiplicativity and hoped-for spectral-sequence or derived-category transitivity statements (pp.~460--461). The proposed Kan--Postnikov complexes use shifted projective resolutions of homotopy groups. To transfer the construction from discrete to schematic objects, he returns explicitly to the cohomology obstruction of p.~444 and hopes to dualize the homology comparison. This is an outline requiring substantial verification, as the author himself stresses. Constant complexes then show why $X_0=e$ is necessary, and the monoid discussion turns the issue into comparison of polynomial with arbitrary cochains. The multiplicative-monoid example on p.~462 is offered as a counterexample to the assumptions previously used; the circled marginal cross is part of the page's evidence and is absent from the arXiv witness.

The response is to ``drop Postnikow and keep Kan'' (p.~462). Since $U(k)$ is not closed under arbitrary fiber products, a schematic horn condition must ensure both representability of the relevant finite limits and a suitably strong surjectivity property (p.~463). The tentative name ``submersions'' covers alternatives ranging from trivial bundles to formally smooth maps admitting sections. The first nontrivial condition asks for the three maps $X_2\to X_1\times X_1$ to be submersions; for a monoid nerve this is meant to force a group. Page 464 lists the desired consequences: the linearization theorem over $\underline Z$, a homology criterion for homotopisms of connected Kan complexes over a general ring, homotopy fibers and Postnikov decomposition, and eventually equivalence of $Hot_n(\underline Z)$ with ordinary $Hot_n$. Page 465 distinguishes restriction of rings from extension, noting that restriction does not preserve projective components without an additional assumption and suggesting resolutions to remove that restriction. Its last sentence proposes compatibility of ring extension with the individual homotopy invariants; the next entry immediately revises that proposal. This transition is essential to reading the section as a developing argument rather than a list of established properties.

\paragraph{113. Soft Postnikov decomposition and a schematic fundamental group (pp. 466--472; 134-7/24--30).}
The entry dated 27.8 corrects the previous evening's assertion. Formula (A) concerns total homology and derived extension of scalars; the analogous termwise assertion (B) for homotopy modules fails without flatness hypotheses (p.~466). The revised arrow in (B) points left, a direction separately emphasized by a circled marginal arrow omitted by the arXiv witness. The reason is the use of shifted projective resolutions in a Postnikov fiber: after tensoring, higher homology can appear, so the base-changed successive fibration need not remain a Postnikov decomposition (p.~467). This explicitly answers the puzzle of p.~444 and motivates the judgment that ordinary Postnikov decomposition is not ``geometric,'' meaning compatible with extension of the ground ring.

The proposed alternative keeps the entire abelianization complex. Using the completed divided-power linearization and the embedding $X\to W(L(X))$, Grothendieck proposes repeatedly taking the homotopy fiber of $X_*\to W(L(X_*))$, without first extracting the first nonzero homology module (pp.~467--468). He names this ``soft Postnikov d\'evissage.'' For a discrete $K(G,1)$ the analogy leads to iterated commutators, and he distinguishes the versions associated with solvability and nilpotence. Formula (D) contains an index problem: the visible inner exponent is $n+1$, like the outer exponent, while the arXiv witness gives $n-1$. The checked text preserves the scan's reinforced plus and does not silently substitute a mathematically expected recurrence. The next page adds the issue of the action of $\pi_1$ on higher homotopy groups; trivial or unipotent action appears relevant to unipotent schematization (p.~469).

The second part of the section revisits complexes with $X_0=e$ but possibly nontrivial $X_1$. A group object gives a fundamental group that need not be abelian, and finite-presentation and smoothness assumptions lead to questions about nilpotence (p.~469). The Kan relations then motivate the presheaf $k'\mapsto\pi_1(X_*(k'))$. The desired invariant is a group object of $U(k)$ representing this sheaf, with an underlying parametrization by a module and a polynomial group law (pp.~470--471). This changes what the notation $\pi_1$ denotes: the group of sections alone is no longer the complete object. Formulas (G) and (G') express evaluation and base-change compatibility for the proposed schematic group. The following duplicated equation on p.~471 is nevertheless literally tautological in the scan, with $\pi_1(X_*(k))$ on both sides; the witness edits one side. The blank between $k$ and $k'$ in the preceding prose likewise contains no arrow on the page.

To admit arbitrary quotient modules, the notion of submersion is reconsidered once more: maps isomorphic to those induced by epimorphisms of modules are proposed as the most convincing of the three possibilities considered so far (p.~471). The printed ``unless \ldots{} does'nt give rise'' contains an apparent extra negation, retained despite the witness's change to ``gives rise.'' The stronger requirement that the schematic fundamental group come from addition on a module is distinguished from commutativity; the discussion ends with tests involving the linearization theorem, a fibration over $K(G,1)$, and the relationship between nilpotent discrete groups and group objects of $U(k)$ (p.~472, referring back to section 94). The handwritten addition about unipotent action on the higher $\pi_i$ belongs here and is not merely incidental: it qualifies the contemplated range of homotopy types. Section 114 begins on the same page by describing the accumulated picture as provisional and preparing a review of the program.

\paragraph{114. A provisional model category and the completion gap (pp. 472--480; 134-7/30--38).}
This review arranges the preceding experiments into a prospective construction. The category $U(k)$ consists of unipotent bundles built from modules; smooth objects correspond to projective modules, submersions to epimorphisms, and smooth submersions to maps with projective kernel (pp.~472--473). The completed divided-power coalgebra describes the unpointed linearization, while the pointed version removes the completion in the basic construction (pp.~474--475). Completion is a substantive distinction: p.~474 explicitly says that the map (C) of p.~467 ``does'nt quite exist'' with $W$ as written, and introduces the completed version $W^{\wedge}$ instead. Base change for the completed functor consequently uses completed tensor products. The subsequent model categories impose $X_i=e$ for $i\leq n$, take sections to pointed semisimplicial sets, and localize at the proposed weak equivalences. The hoped-for comparison with ordinary homotopy types is restricted to the appropriate connected range over $\underline Z$ (p.~476). On p.~477 the first expression for the category of models visibly uses $\pi_n(k)$, although the surrounding formulas use the handwritten model-category symbol $\mathcal M_n(k)$; this notation anomaly is retained rather than silently regularized.

Smooth Kan resolutions are then proposed to derive extension of scalars and the total linearization $LH.(X_*)=L_{pt}(X_*)$ (p.~477). The next two pages ask whether this homology functor is conservative and whether pointed linearization and completed linearization are quasi-isomorphic. They also reconsider comparison with ordinary free-module homology. The second question undergoes an immediate correction: on p.~479 the author says the proposed comparison with $k^{(X_*(k))}$ is generally false away from $k=\underline Z$. In the first nonvanishing degree, the expected $k$-module invariant is identified with integral homology of the section complex; tensoring that homology with $k$ produces a different candidate. The word \emph{integral} is explicit in the scan, whereas the arXiv witness gives \emph{internal}. The correction renews the need to construct homotopy modules intrinsically and to obtain them through a Postnikov argument, rather than transfer all ordinary homology assertions indiscriminately.

The final page makes the outstanding construction precise. The exponential map gives $X_*\to W^{\wedge}(L(X_*))$, but the required homotopy-class map is $X_*\to W(L_{pt}(X_*))$ (p.~480). Even a quasi-isomorphism between the two linearizations does not by itself remove the difference between $W^{\wedge}$ and $W$; the handwritten phrase ``flaw in this `argument''' marks exactly this gap. The target of the completed map is a semisimplicial sheaf and need not belong to the proposed model category. Hard or soft Postnikov decomposition therefore still requires this factorization, a suitable notion of homotopy fiber, and homotopy constructions internal to the model category. The failure of $U(k)$ to have fiber products explains why Quillen's standard machinery cannot simply be applied as stated. These are explicit unresolved tasks, and section115 returns the following morning to the meaning of the coefficient sheaf itself.

\paragraph{115. Quasicoherent coefficients and prorepresentable linearization (pp. 481--489; 134-7/39--47).}
The entry dated28.8 restarts the formula numbering and identifies the source of the preceding confusion. The appropriate free linearization is over the quasicoherent ring sheaf $\underline O_k$ on the fpqc site, not over a constant ring (pp.~481--482). Grothendieck explains that the obstruction groups in schematic Postnikov decomposition use the quasicoherent sheaf associated with a homotopy module, rather than the same module regarded as a constant coefficient group. The circled marginal clarification of the argument $X(n)_*,\pi_{n+1}$ belongs to this distinction; the overwritten version in the running formula remains marked as doubtful. He connects the earlier mistake to having translated geometric language into commutative algebra too thoroughly, losing sight of which coefficient objects were intended.

The universal free Module $\underline O_k^{(X)}$ gives the Hom, Ext and derived-Hom expressions (4)--(5') (p.~483), but is generally not quasicoherent. The proposed replacement would represent the same functor on quasicoherent coefficients. For $X=W(M)$, the continuous-Hom calculation through truncated divided-power modules shows that the correct object is a proobject, rather than the completed module regarded as an ordinary module (pp.~484--485). The obstruction is again scalar extension: inverse limits do not generally commute with tensoring by a new ground ring. Formula(13) packages the finite products of divided powers into a strict projective system; ordinary representability would require it to become essentially constant. The added phrase ``but with slightly different meaning'' therefore changes the interpretation of $L$ in (14), now valued in pro-modules. The pointed functor still takes values in ordinary modules, embedded as constant proobjects. Formula(18), on p.~486, supplies the exponential map into $WL(X)$ in the procategory and distinguishes this abelianization from its richer linear structure.

The proposed derived comparison (20) and its Ext form are then tested on $X=W(M)$ (pp.~486--487). Formula(21) contains a meaningful handwritten revision from a limit expression to a direct sum of Ext groups of divided powers. Projectivity of $M$ gives the expected vanishing in positive degree, with a passage to finite-type submodules sketched for the infinite case. For nonprojective $M$, the potential nonzero Ext$^1$ term challenges the author's previous expectation that quasicoherent cohomology of these bundles should vanish. His reference to Larry Breen is a question about this expectation, not a proof of the general formula. The text visibly says ``quasiunipotent'' where the witness shortens it to ``unipotent'' (p.~487). The next page restricts attention to the smooth case and postpones the projective-resolution argument for full generality.

For a semisimplicial complex, (22) uses Hyperext rather than a degreewise Ext construction; (23) compares the resulting pro-linearization with pointed linearization (p.~488). The comparison map(24) revisits the map from p.~454 with its proobject interpretation made explicit. Its relationship to a quasi-isomorphism, and then to the componentwise inverse limit used on p.~478, still requires foundational work. Page489 plainly admits that no argument has yet been found for the desired assertion under the proposed smoothness, connectivity and possibly Kan hypotheses. The second comparison is reformulated as $\underline O_k^{(X_*)}\to L(X_*)$ in (25). The author does not confidently assert that this is a quasi-isomorphism: quasicoherence of the first homology sheaves and essential constancy of the second remain in doubt. The practical hope is instead that it gives the needed comparisons against quasicoherent coefficients through(20). Section116 subsequently reports the collapse of the earlier general abelianization theorem, so these provisional claims must be read with that later revision in view.

\paragraph{116. The failed abelianization theorem and the interruption of the notes (pp. 490--491; 134-7/48--49).}
The entry dated 7.9 reports the decisive revision of the preceding program. After finding a simple description of schematic homotopy groups, Grothendieck learned in a telephone conversation with Illusie that the asserted comparison between discrete and schematic homology or cohomology was false. The anticipated equivalence between schematic and ordinary homotopy types therefore fails (p.~490). This is an explicit retraction of a central assumption underlying sections111--115; it does not cause him to abandon schematic homotopy types, whose formal properties he now wants to study in their own right, possibly through different models. The remainder of the section explains the intervening pause in writing through fatigue and the need to attend to the death of his granddaughter Ella. Page491 distinguishes work allowed to take place through attention from a compulsive forward movement that leaves perceptions disconnected from one's picture of reality. The revisions matter to the voice of this passage: the hand moves ``who'' after ``finally,'' retains ``his language'' for the body's intervention, and changes ``cause'' to ``caused.'' The witness substitutes ``which'' and ``its'' and loses the last correction. These pages provide both the mathematical warning needed to read the preceding conjectures and the author's own account of why the notebook paused before the technical review resumes in section117.

\paragraph{117. Size, exactness, and the distinction between schemes and formal schemes (pp. 492--496; 134-7/50--54).}
The entry dated 10.9 begins with two technical qualifications. The affine site used earlier is not automatically small in the working universe; a possible remedy is to bound the sizes of the modules and the rings used as arguments, using the fppf site for finite-type bundles (p.~492). More consequentially, $M\mapsto W(M)$ is right exact but need not preserve kernels. Thus a free resolution of $M$ need not become a resolution of $W(M)$. The derived-Hom comparison used to interpret formula(21) of p.~487 is therefore questioned in general and retained for the flat case (p.~493). This is a specific dependency of the cohomology calculation, not merely a set-theoretic caveat.

The finite projective case clarifies the geometry through duality. Polynomial maps can be described by maps of symmetric algebras or by continuous maps of the dual completed coalgebras; preserving the origin translates into augmentation compatibility (pp.~493--494). The uncompleted divided-power coalgebra is instead dual to the adically completed symmetric algebra. Consequently, the two algebraic descriptions lead to schemes and formal completions along their zero sections, respectively. The scan's formulas require care: the prose on p.~493 includes a dual check on $M$, but formulas(2), (3) and (4) omit dual signs that the arXiv witness supplies. Those omissions remain in the checked text. The upward-arrow-shaped completion mark in(3') on p.~494 is likewise recorded as drawn, and a separate unreadable pen mark above the star in(5) remains marked.

This geometric interpretation overturns the earlier hoped-for equivalence between the pointed and nonpointed linearizations (pp.~495--496). For the additive group in characteristic $p$, the degree-one comparison distinguishes additive polynomials from formal series $\sum_{i\geq0}c_i t^{p^i}$: the schematic side allows only finitely many nonzero coefficients. Grothendieck acknowledges that this example is neither simply connected nor over $\underline Z$, but states that he also no longer expects the comparison to be a quasi-isomorphism for the basic connected examples previously contemplated. He reports that Illusie rejected the vanishing assumptions behind his attempted ``way-out'' argument. The closing underline therefore warns explicitly against substituting $L_{pt}(W(M))$ for $L(W(M))$ when computing schematic homology or cohomology. Formal invariants have their own meaning and are distinct. Section118 begins on the same page with a different and deliberately striking claim: homotopy invariants themselves should nevertheless be recoverable from the first-order tangent object.

\paragraph{118. The tangent complex as a proposed homotopy invariant (pp. 496--500; 134-7/54--58).}
Against the distinction just established between schematic and formal cohomology, Grothendieck proposes that homotopy groups should depend only on the first-order tangent object. The Lie functor sends a smooth pointed bundle $W(M)$ to $M$, and formula(9) identifies the homotopy groups of $X_*(k)$ with those of the semisimplicial module $\mathrm{Lie}(X_*)$, hence with homology of its associated chain complex (p.~496). The next page sharply qualifies the claim: he has no direct description of a natural isomorphism under the stated hypotheses. It is tautological for complexes arising from chain complexes and is expected to extend through Postnikov decomposition, if every relevant model can be replaced by one admitting such a decomposition; naturality would still require work (p.~497). The typed revision makes this formula a ``cornerstone'' of a theory of schematic homotopy types and a proposed intrinsic definition of the invariants, rather than subordinating them from the outset to discrete homotopy groups. Weak equivalences would then be detected by the tangent complexes. Both subsequent comparison displays map an $X_*$ object to a $Y_*$ object; the arXiv witness repeats $X_*$ as target in each.

The contemplated application is to explicit models of spheres (pp.~498--499). To make the Lie formula useful for ordinary homotopy groups, one must first realize the discrete type schematically over $\underline Z$ and obtain a workable model. Successive suspension from $S^1$ suggests the need for suspension constructions, complementing the earlier need for fibers and loops; the proposed derivator would organize the finite integration and cointegration operations. These are motivations and difficult construction problems, not completed calculations of sphere homotopy groups. The discussion of $S^2$ as a homogeneous space under $S^3$, and of the corresponding group/flag-variety geometry, asks for a model of the entire relationship: a group, a subgroup playing the role of a torus or Borel subgroup, and a quotient. Several overprinted group-name letters remain uncertain in the text, while the star on $\underline C^*$ is visibly an asterisk rather than the witness's multiplication cross.

Hemispherical complexes supply a second route (pp.~499--500). Their smaller list of boundary and degeneracy operations, together with a prospective direct computation of homotopy in the discrete setting, is compared with the tangent-complex method. The marginal exclamation and vertical rule emphasize this alternative. The author immediately admits that he has not checked the proposed computation for arbitrary hemispherical sets that do not arise from an infinity-groupoid or infinity-stack. Comparing the two descriptions is therefore offered as a possible route toward understanding formula(9), with that unresolved premise preserved. Section119 then resumes on26.9 after another interruption and returns to the failure of the earlier homology comparison.
